\nomenclature[A]{IA}{Inteligência Artificial}
\nomenclature[A]{OAE}{Obras de Artes Especiais}
\nomenclature[A]{CV}{Visão Computacional}

\chapter{Introdução}
\label{chap:intro}.
O monitoramento de progresso é um dos elos fundamentais para gestão de projetos, auxiliando na conclusão bem sucedida da construção de edificações e infraestruturas (Pal et al., 2023). Na construção de Obras de Arte Especiais (OAE), como as pontes, as condições diferem significativamente das edificações convencionais, caracterizando-se por processos altamente setorizados, distribuição espacial discreta e uso intensivo de componentes pré-fabricados (Zhang, 2023). Embora essas estruturas apresentem padrões geométricos repetitivos e similares, a extensão do canteiro e a complexidade logística impõem desafios severos à gestão e ao acompanhamento físico do projeto (Zhang, 2023).

Os métodos tradicionais de monitoramento do progresso se baseiam em registros manuais, inspeções \textit{in loco} e atualizações em planilhas ou papel (Pal et al., 2023). Nesse sentido, essa abordagem apresenta desvantagens críticas: é excessivamente lenta, trabalhosa, propensa a erros humanos e influenciada pela subjetividade dos inspetores (Qureshi et al., 2023). Devido à incapacidade de abranger a totalidade do canteiro com alta frequência, os dados coletados ficam rapidamente desatualizados, tornando a tomada de decisão reativa e propiciando atrasos e estouros de orçamento (Reja et al., 2022). Ademais, a simples captura visual do canteiro sem processamento automatizado não soluciona o problema, pois apenas transfere o gargalo operacional do campo para o escritório.  



 Nesse cenário, a integração de ferramentas como a Modelagem da Informação da Construção (BIM) a técnicas de Visão Computacional (CV) em Inteligência Artificia (IA)l tem se consolidado como a via mais eficiente para o monitoramento automatizado (Pal et al., 2023). O uso de \textit{Deep Learning} permite processar dados visuais complexos, como nuvens de pontos, comparando o modelo planejado \textit{as-planned} com a execução real \textit{as-built} para extrair informações objetivas de avanço (Reja et al., 2022; Braun et al., 2020).  
 FALTOU FALAR E CITAR BENEFICIOS DA FOTOGRAMETRIA. 
 
 Apesar desses avanços, a literatura e as aplicações de monitoramento automatizado ainda concentram-se majoritariamente em edificações verticais, revelando uma lacuna no desenvolvimento de fluxos adaptados às especificidades geométricas e operacionais das pontes. Portanto, este trabalho propõe um fluxo semi-automatizado para o monitoramento de pontes, integrando fotogrametria, eliminação de ruídos de nuvens de pontos por \textit{deep learning} (RandLA-Net) e algoritmos de alinhamento geométrico (OBB-ICP) ao modelo BIM 4D, visando eliminar a carga manual, fornecer dados acionáveis aos gestores e mitigar ineficiências construtivas.

\section{Objetivos}
\label{sec:objetivos}
O objetivo geral deste trabalho é propor (um sistema de monitoramento e progressão de obras usando inteligência artificial, para gestão de obras de infraestrutura - ou ponte, visando...)  

Para alcançar este propósito, definem-se os seguintes objetivos específicos:

\begin{itemize}
	\item Mapear a captura de dados, estabelecer o protocolo de aquisição de imagens fotogramétricas (?), processamento e geração de nuvens de pontos densas de obras de pontes (as-built), sintetização de nuvem de pontos e extração de informações do modelo BIM, talvez fazer uma discussão também sobre o georreferenciamento
	
	\item Aplicar Visão Computacional / Deep Learning (RANDLANET) para implementar/adaptar algoritmos de IA para segmentar (ou remover ruído) e a segmentação dos elementos estruturais da ponte a partir das nuvens de pontos. Todas as etapas da metodologia: Data collection, preparation…
	
	\item Executar a verificação comparativa: utilizar algoritmos (nesse caso falo do obbp-icp) para integrar e alinhar a nuvem de pontos processada ao modelo 3D/4D BIM para identificar avanço, atrasos ou discrepâncias de execução.
	\item Mensurar a acurácia geométrica do alinhamento, o tempo computacional e o esforço operacional do método proposto, analisando sua aplicabilidade prática frente aos processos tradicionais de medição visual e manual.
	
\end{itemize}

\section{Organização do Trabalho}
\label{sec:organizacao}
Este trabalho está estruturado em seis capítulos, organizados da seguinte forma:

O Capítulo \ref{chap:intro} apresenta a introdução e contextualização do tema, a justificativa para a realização da pesquisa, bem como os objetivos propostos.

O Capítulo \ref{chap:revisao} dedica-se à revisão da literatura, abordando conceitos fundamentais sobre

O Capítulo \ref{chap:metodologia} descreve a metodologia da pesquisa, detalhando as etapas de construção do modelo conceitual, a modelagem da 

O Capítulo \ref{chap:resultados} expõe o desenvolvimento e validação do modelo,

O Capítulo \ref{chap:discussao} faz a discussão sobre o resultado da 

Por fim, o Capítulo \ref{chap:conclusao} apresenta as considerações finais, as limitações encontradas ao longo do estudo e as sugestões para trabalhos futuros.

\section{Contexto da Pesquisa}
A pesquisa desenvolvida neste Trabalho de Conclusão de Curso está inserida nas atividades do Grupo de Pesquisa e Extensão em Gestão e Tecnologia das Construções (GETEC), vinculado à Escola Politécnica da Universidade Federal da Bahia (UFBA).

O trabalho foi realizado no contexto de um projeto de pesquisa em nível de doutorado, conduzido por \citeonline{tesecarol}, com foco no desenvolvimento de um sistema de Gêmeo Digital para gerenciamento e controle de recursos físicos em sistemas construtivos de parede de concreto. O presente estudo se integra como contribuição complementar, sob orientações do Professor Carlos Eduarno Viana Nunes. Durante o desenvolvimento desta pesquisa, eu atuei como bolsista de Iniciação Tecnológica e Industrial (ITI),sob orientação da Professora Dayana Bastos Costa e em colaboração com a Universidade de São Paulo (USP), representada pelo Professor Fabiano Rogerio Corrêa. Minha participação deu continuidade no trabalho iniciado pelo bolsista Caio Vinicios Borges de Souza, onde atuei auxiliando na pesquisa e na produção do trabalho, com fomento do Conselho Nacional de Desenvolvimento Científico e Tecnológico (CNPq).

A “Empresa A” foi selecionada como parceira deste estudo. Trata-se de uma empresa brasileira de médio porte do setor da construção civil, atuante no estado da Bahia desde 1987 e parceira do GETEC desde 2019. A empresa participou do projeto no contexto do programa MAI/DAI do CNPq, incentivando iniciativas de inovação tecnológica e melhoria dos processos de gestão. O canteiro de obras analisado pertence a um empreendimento residencial vinculado ao programa habitacional "Minha Casa, Minha Vida", servindo como base para as observações, coleta de informações e desenvolvimento das pesquisas adotadas neste trabalho.

O uso de ferramentas de Inteligência Artificial Generativa utilizadas neste trabalho ocorreu em conformidade às normas do Guia para Uso Ético e
Responsável da Inteligência Artificial Generativa na
Universidade Federal da Bahia, disponível em \cite{ufba2025ia}.

